\section{Introduction}
\label{sec:intro}

Multilingual deployment of large language models (LLMs) exposes a critical calibration gap: non-English users receive substantially worse confidence estimates, with Expected Calibration Error (ECE) up to 5$\times$ higher in low-resource languages \cite{zhou2025lace}. Standard post-hoc methods such as temperature scaling \cite{guo2017calibration} require per-language validation data, and verbalized confidence \cite{kadavath2022language} operates on output logits without accessing internal representations. A recent line of work identifies an alternative signal: LLM hidden states encode a ``correctness direction'' nearly orthogonal to verbalized confidence (cosine similarity below 0.04) \cite{miao2026closing}, and correctness probes trained on this direction achieve approximately 50\% ECE improvement over output-based baselines \cite{miao2026coral}. Both CCPS \cite{khanmohammadi2025ccps} and the Confidence Manifold framework \cite{confidence_manifold2026} confirm that calibration-relevant information concentrates in low-dimensional representational subspaces. All of these probe-based calibration methods, however, have been developed and evaluated exclusively in English.

No prior work has tested whether correctness probes transfer across languages. The closest studies address related but distinct targets. \citet{xiao2025knowledge} probe knowledge boundaries (an input-level property of whether a question is answerable), not correctness (an output-dependent property of whether the model answered correctly). \citet{civelli2026shared} probe item difficulty using Item Response Theory (IRT) labels, a language-agnostic question property independent of any particular model's response. Neither study reports calibration metrics (ECE, Brier score) or compares multiple transfer strategies. DALI \cite{ravisankar2026dali} provides causal evidence through activation patching that English middle-layer representations can repair errors in 24 non-English languages, which confirms that transferable correctness information exists but does not quantify when or why probe-level transfer fails. This gap motivates three research questions: (1) Does the correctness direction persist across languages, or is it language-specific? (2) How does cross-lingual transfer depend on representation geometry, specifically layer depth and cross-lingual representational similarity? (3) Can distributional matching shift the transfer frontier without target-language correctness labels?

We present the first large-scale study of correctness probe transfer across architectures. The study covers 2+ model families, 20 languages in three resource tiers, and 5 transfer methods evaluated on GlobalMMLU. In 7B-scale pilot experiments, deeper layers generally improve both English probe quality and cross-lingual transfer (Spearman $\rho > 0.7$, $p < 0.05$ for 4 of 5 languages), but the Pareto frontier reveals a local tradeoff at the deepest layers (Figure~\ref{fig:pareto}): the final 1--3 layers gain marginal English AUROC (+0.8\,pp) at the cost of transfer degradation (1.5--4\,pp in 3 of 5 languages). The optimal operating point is language-dependent, and the local tradeoff shape varies across architectures: Qwen exhibits consistent deep-layer degradation, while Llama displays language-specific patterns.

This tradeoff conceals a deeper geometric structure. The raw correlation between Centered Kernel Alignment (CKA) \cite{kornblith2019cka} and transfer Area Under the ROC Curve (AUROC) across all language-layer pairs is negative, because both quantities covary with depth. After controlling for depth via partial correlation, the sign reverses: within each layer, higher cross-lingual CKA predicts better probe transfer. In 7B-scale pilot analysis, this partial correlation is significant for Llama-3.1-8B ($\rho = +0.41$, $p = 0.006$); Qwen2.5-7B shows a consistent but underpowered trend ($\rho = +0.18$, $p = 0.274$, $n=40$). Full-scale validation across 20 languages (\S\ref{sec:geometric}) will determine generality. This Simpson's Paradox exposes a confounding pattern that naive cross-layer pooling obscures: what appears as a negative CKA-transfer association in aggregate is in fact a positive within-layer relationship.

A methodological finding further qualifies these results. We apply the control task framework of \citet{hewitt2019control} to cross-lingual probing. In 7B-scale pilot experiments, English probing controls reveal that selectivity concentrates in deeper layers (L20+, Hewitt \& Liang selectivity $> 0.20$), raising concerns that shallow-layer positive transfer in low-resource languages may reflect capacity artifacts rather than genuine linguistic signal; cross-lingual control experiments are planned to quantify this effect (\S\ref{sec:controls}). This is the first application of probing controls to cross-lingual transfer evaluation. On the practical side, linear projection via Procrustes alignment achieves near-optimal transfer with approximately 50 parallel sentence pairs, though saturation rate depends on resource tier. We also design Calibration-aware Cross-Lingual Regularization (CaCLR), a distributional matching objective ($\mathcal{L}_{\text{BCE}} + \alpha \cdot \mathcal{L}_{\text{SB-ECE}} + \beta \cdot \mathcal{L}_{\text{consistency}}$), as a principled test of whether Maximum Mean Discrepancy (MMD) alignment can improve zero-correctness-label transfer (results pending experimental validation).

Our experimental framework provides systematic coverage along four axes, with full-scale replication on frontier-scale models. We test 2+ architectures (Qwen3-235B-A22B and Llama 4 Scout 17B-16E as primary models) across 20 languages that span high, medium, and low resource tiers. We compare 5 transfer methods ordered by complexity: zero-shot application, mean-shift alignment, Procrustes rotation, few-shot probe retraining, and CaCLR (\S\ref{sec:transfer}). All transfer methods except few-shot retraining satisfy a zero-correctness-label constraint and need no target-language correctness annotations. We decompose probe targets into difficulty-shared and capability-specific components through IRT-based projection (\S\ref{sec:decomposition}), which explains why transfer succeeds at some layers and fails at others. CKA geometric analysis with depth-controlled partial correlation provides the diagnostic framework to explain failure modes (\S\ref{sec:geometric}).

Our contributions are:
\begin{enumerate}
\item \textbf{Architecture-dependent Pareto frontier and CKA geometric analysis.} We discover that deeper layers generally improve both English probe quality and cross-lingual transfer, but the Pareto frontier reveals a local tradeoff at the deepest layers whose shape varies across architectures. A Simpson's Paradox in the CKA-transfer relationship reveals that within-layer representational similarity predicts transfer success, a confounding pattern obscured by naive cross-layer pooling (\S\ref{sec:geometric}).

\item \textbf{Difficulty-capability decomposition.} We decompose cross-lingual transfer into a difficulty-shared component (transferable at all depths) and a capability-specific component (responsible for deep-layer degradation), which yields prescriptive predictions for alignment strategy selection (\S\ref{sec:decomposition}; experiments in progress).

\item \textbf{CaCLR distributional matching test.} We design CaCLR ($\mathcal{L}_{\text{BCE}} + \alpha \cdot \mathcal{L}_{\text{SB-ECE}} + \beta \cdot \mathcal{L}_{\text{consistency}}$) as a principled baseline that tests whether MMD-based distributional matching can improve zero-correctness-label transfer. Results are pending experimental validation (\S\ref{sec:transfer}).

\item \textbf{Empirical contributions with practical implications.} (a) Control task analysis reveals that probe selectivity concentrates in deeper layers, raising concerns about capacity artifacts in low-resource shallow-layer transfer, a field-level methodological warning (\S\ref{sec:controls}). (b) Linear projection with approximately 50 parallel pairs achieves near-optimal transfer with resource-tier-dependent saturation. (c) Procrustes alignment exhibits effect reversal: it helps when zero-shot AUROC $< 0.50$ but degrades performance when zero-shot AUROC $> 0.60$ (\S\ref{sec:transfer}).
\end{enumerate}
